\section{Ordered cache experiments across three models}
\label{sec:multimodel-cache-100}
\label{sec:oss-cache-100}

The completed cache study compares GPT-OSS, Qwen and Gemma on the same
100 held-out source problems, with one baseline and one fixed positive
bias-$+1$ generation per model and problem. Source-problem identities
match across all three cohorts. The deterministic selection contains
9 AIME24, 9 AIME25, 12 AMC23, 24 MATH-500, 23 Minerva and 23 OlympiadBench
problems; smaller held-out pools are exhausted before the larger pools.
Each model retains its own frozen calibration policy. Per-model paired
prompts, seeds and generation settings match. Generation uses temperature
0.6, top-$p$ 0.95, a 32,768-token output limit, and synchronous
single-sequence execution. These are separate generations from the
1,395-attempt detailed runs and do not constitute independent seed
replications of those runs.

\paragraph{Measurement and cache controls.}
Ordered router-hook top-$k$ sets cover all 24 OSS, 40 Qwen and 30 Gemma
layers. All prefill chunks are excluded, so $N$ returned tokens correspond
to $N-1$ recorded decode steps at every layer. Source hashes, paired inputs,
layer coverage and trace lengths are validated. Hook top-$k$ is not an
independent observation of native dispatch; tied scores may differ.

Each answer starts with a cold per-layer cache. Initial pinned-expert loads
count, even if a pin is never subsequently used. The four conditions are
baseline LRU, baseline with calibration-frequency pins, guided LRU, and
guided routing with target pins; remaining slots use LRU. Pins are restricted
to each model's policy layers, with equal pin counts for frequency and target
controls. Other layers use LRU. Simultaneous top-$k$ requests cannot evict
one another; sorted expert IDs break simultaneous recency ties.

OSS uses four of 32 experts and budgets of 8, 16 and 32 slots per layer.
Qwen uses eight of 256 and Gemma eight of 128, with budgets of 16, 32 and
64 slots. Capacities are equal across conditions within a model. They do
not represent equal bytes or equal fractions of expert capacity across
models. Shared experts, dense operations and KV-cache memory are outside
the expert-load metric. No SSD reads, transfer bandwidth or latency are measured.

\paragraph{Accuracy and output length.}
Offline extracted-answer equivalence rescoring gives baseline/guided
correct counts of 76/75 for OSS, 83/82 for Qwen, and 84/84 for Gemma.
Their accuracy changes have 95\% intervals of $[-7,+5]$, $[-5,+3]$ and
$[-4,+4]$ percentage points, respectively. These use 5,000 paired,
dataset-stratified source-problem resamples with seed 42. The original
saved scorer instead gives 53/48, 62/61 and 59/59; those labels remain
unchanged in the source artifacts. No model establishes accuracy
non-inferiority from these intervals alone.

OSS outputs are $22.04\%$ longer on average; Qwen outputs are $1.06\%$
shorter and Gemma outputs $1.61\%$ shorter. Token-limit counts are 2/6,
7/7 and 3/5, respectively. The distinct length responses make total loads
per answer necessary alongside loads per token.

\begin{table}[ht]
\centering\small
\caption{Matched 100-problem cache experiments. B/G denotes baseline/guided. Accuracy uses offline equivalence rescoring; changes and intervals are percentage points. Token changes are relative percentages.}
\label{tab:cache-multimodel-quality}
\resizebox{\linewidth}{!}{%
\begin{tabular}{llrrrr}
\toprule
Model & Correct B/G & Change [95\% CI] & Mean tokens B/G & Token change (\%) & Capped B/G \\
\midrule
GPT-OSS & 76/75 & -1.0 [-7.0, +5.0] & 4,418.9/5,392.9 & +22.04 & 2/6 \\
Qwen & 83/82 & -1.0 [-5.0, +3.0] & 12,116.1/11,988.0 & -1.06 & 7/7 \\
Gemma & 84/84 & +0.0 [-4.0, +4.0] & 9,645.4/9,490.5 & -1.61 & 3/5 \\
\bottomrule
\end{tabular}
}
\end{table}
\begin{table}[ht]
\centering\small
\caption{Simulated mean expert loads per complete answer, including cold-start pin loads, across all routed layers. Capacity is slots per layer; frequency and target pins use the same slot budget within each model.}
\label{tab:cache-multimodel-loads}
\resizebox{\linewidth}{!}{%
\begin{tabular}{llrrrr}
\toprule
Model & Slots & Baseline LRU & Baseline frequency & Guided LRU & Guided targets \\
\midrule
GPT-OSS & 8 & 134,007 & 134,747 & 158,903 & 162,649 \\
GPT-OSS & 16 & 40,080 & 40,068 & 45,581 & 46,092 \\
GPT-OSS & 32 & 723 & 723 & 717 & 717 \\
Qwen & 16 & 2,400,863 & 2,410,567 & 2,266,939 & 2,252,038 \\
Qwen & 32 & 1,900,161 & 1,901,191 & 1,772,502 & 1,767,504 \\
Qwen & 64 & 1,184,307 & 1,184,273 & 1,092,516 & 1,091,858 \\
Gemma & 16 & 827,247 & 821,327 & 825,811 & 912,591 \\
Gemma & 32 & 354,596 & 352,565 & 353,664 & 378,813 \\
Gemma & 64 & 60,310 & 60,284 & 59,939 & 61,556 \\
\bottomrule
\end{tabular}
}
\end{table}
\begin{table}[ht]
\centering\small
\caption{Relative change (percent) for guided target pinning versus baseline LRU, with paired dataset-stratified problem-bootstrap 95\% intervals. Negative values mean fewer simulated loads. Intervals use 5,000 resamples, seed 42, and no multiplicity correction.}
\label{tab:cache-multimodel-effects}
\resizebox{\linewidth}{!}{%
\begin{tabular}{llrr}
\toprule
Model & Slots & Loads/answer: change [CI] & Loads/token: change [CI] \\
\midrule
GPT-OSS & 8 & +21.37 [-8.92, +62.53] & -0.55 [-2.74, +1.66] \\
GPT-OSS & 16 & +15.00 [-13.33, +53.22] & -5.77 [-9.06, -2.56] \\
GPT-OSS & 32 & -0.78 [-1.05, -0.50] & -18.71 [-39.53, +8.83] \\
Qwen & 16 & -6.20 [-16.35, +5.50] & -5.20 [-6.34, -4.00] \\
Qwen & 32 & -6.98 [-17.12, +4.72] & -5.99 [-7.19, -4.65] \\
Qwen & 64 & -7.81 [-14.90, +0.08] & -6.82 [-10.26, -2.14] \\
Gemma & 16 & +10.32 [+2.69, +18.68] & +12.12 [+10.23, +14.44] \\
Gemma & 32 & +6.83 [-0.72, +15.07] & +8.57 [+5.84, +11.73] \\
Gemma & 64 & +2.07 [-4.74, +9.97] & +3.73 [+0.08, +7.76] \\
\bottomrule
\end{tabular}
}
\end{table}


\paragraph{Qwen: a promising loading reduction.}
Guided target pinning reduces mean simulated loads per answer by
$6.20\%$, $6.98\%$ and $7.81\%$ at 16, 32 and 64 slots per layer.
It also uses fewer loads than the same-budget frequency-pinned baseline.
Guided LRU already achieves most of the reduction; target pinning adds a
smaller improvement on the same guided traces. Thus the result belongs
primarily to the changed routing trajectories, with a smaller incremental
benefit from reserving slots for the targets. The accuracy point estimate
is one fewer correct answer out of 100, and output length is slightly lower.
Per-token reductions are $5.20\%$, $5.99\%$ and $6.82\%$, with
all three 95\% intervals below zero. All per-answer intervals still
include zero: $[-16.35,+5.50]\%$, $[-17.12,+4.72]\%$ and
$[-14.90,+0.08]\%$. Thus per-token locality improves under the reported
resampling criterion, while a population-level reduction in total loads
per answer remains unresolved. These are not end-to-end speed measurements.

\paragraph{Gemma: pinning rare targets is costly.}
Guided target pinning increases loads per answer by $10.32\%$, $6.83\%$
and $2.07\%$ at 16, 32 and 64 slots. At 16 slots the 95\% interval is
$[+2.69,+18.68]\%$; at 32 and 64 slots the intervals include zero.
Loads per token increase at all three budgets. Ordinary LRU on the guided
traces instead gives 825,810.8, 353,664.2 and 59,938.7 loads per answer,
slightly below the respective baseline-LRU means. Explicit target pinning
therefore consumes capacity that adaptive LRU uses more effectively.
This agrees with the full-diagnostics finding that the accuracy-associated
Gemma targets remain infrequent even after bias. There is no aggregate
accuracy difference in the 100-problem sample.

\paragraph{OSS: local savings are outweighed by longer answers.}
At 16 slots, guided target pinning reduces loads per token by $5.77\%$
($[-9.06,-2.56]\%$), but increases observed loads per answer by $15.00\%$
($[-13.33,+53.22]\%$). At eight slots, the per-answer increase is
$21.37\%$ ($[-8.92,+62.53]\%$). Ordinary LRU uses fewer loads than target
pinning on the same guided traces. At 32 slots all OSS experts fit, so the
small reduction in loads concerns initial loading of distinct experts,
not reduced eviction under memory pressure.

\paragraph{Interpretation.}
The three models do not support a universal cache benefit from promoting
accuracy-associated identities. Qwen is the strongest candidate for a
subsequent offloaded-inference benchmark; Gemma exposes the cost of pinning
rare targets, and OSS shows how longer outputs can offset per-token locality.
A small accuracy loss could be acceptable in exchange for sufficient speed,
but establishing that trade-off requires actual answer latency under equal
memory and transfer budgets. Predictive prefetching is not evaluated here.

Cache intervals resample paired problems within datasets (5,000 resamples,
Python random seed 42), recomputing ratios of total loads and token counts.
They condition on the frozen policies and saved generations and are not
adjusted for multiple budgets or models. They do not include independent
seed or policy-fitting variation. Prefill warming, cross-answer reuse,
batching, prefetch and transfer costs are excluded. Weights remain resident
during generation; instrumentation time cannot be interpreted as SSD-backed
inference time. The matched cohort strengthens descriptive comparison but
does not isolate architectural causes of differences between models.
